\documentclass[11pt]{article}
\usepackage{amsmath,amssymb,amsthm}
\usepackage[margin=1in]{geometry}
\usepackage[T1]{fontenc}
\usepackage{lmodern}

\newtheorem{theorem}{Theorem}
\newtheorem{lemma}[theorem]{Lemma}
\theoremstyle{definition}
\newtheorem{definition}[theorem]{Definition}

\title{The bridge implication for the \((3,4,\infty)\) lattice form\\
\large Restriction lemma for \(\Omega_0\)}
\author{Benjamin Stanley Frohman}
\date{4 October 2026}

\begin{document}
\maketitle

\begin{abstract}
This note records the lattice bridge already checked in
\texttt{OmegaZero34/TetherInterface.lean}.
The classification is \texttt{uniqueness\_over\_Q}.
The bridge implication is \texttt{restriction\_lemma}:
nondegeneracy forces the scalar to be nonzero.
It does not prove the geometric hypotheses, and it does not prove \(X\simeq S^6\).
\end{abstract}

\section{The form}
On the rank-\(4\) lattice with basis \((\gamma,u,w,\delta)\),
\[
\Omega_0=
\begin{pmatrix}
0&0&0&1\\
0&0&6&0\\
0&-6&0&0\\
-1&0&0&0
\end{pmatrix}.
\]
Indices below are \(0\)-based, so \(\Omega_{0,03}=1\) and the scalar in the Lean statement is the \((0,3)\)-entry.

\section{Statements}

\begin{definition}[Invariant form]
An invariant rational form on this representation is a matrix
\(\omega\in\mathrm{Mat}_4(\mathbb{Q})\) satisfying
\begin{align*}
\omega^{\mathsf T}&=-\omega,\\
T_1^{\mathsf T}\omega T_1&=\omega,\\
T_2^{\mathsf T}\omega T_2&=\omega,\\
\det\omega&\neq 0.
\end{align*}
In the Lean development this is the structure \texttt{InvariantFormQ}.
\end{definition}

\begin{theorem}[Uniqueness over \(\mathbb{Q}\)]
Let \(\Psi\in\mathrm{Mat}_4(\mathbb{Q})\) satisfy
\(\Psi^{\mathsf T}=-\Psi\),
\(T_1^{\mathsf T}\Psi T_1=\Psi\), and
\(T_2^{\mathsf T}\Psi T_2=\Psi\).
Then
\[
\Psi=\Psi_{03}\,\Omega_0.
\]
The scalar \(\Psi_{03}\) may be zero.
This is \texttt{uniqueness\_over\_Q}.
\end{theorem}

\begin{theorem}[Bridge implication]\label{thm:bridge}
Let \(\omega\in\mathrm{Mat}_4(\mathbb{Q})\) be alternating, preserved by \(T_1\) and \(T_2\), and nondegenerate.
Then
\[
\omega=\omega_{03}\,\Omega_0
\qquad\text{and}\qquad
\omega_{03}\neq 0.
\]
Equivalently, \(\omega=\lambda\Omega_0\) for \(\lambda=\omega_{03}\in\mathbb{Q}^{\times}\).
This is \texttt{restriction\_lemma}.
\end{theorem}

\begin{proof}
By uniqueness over \(\mathbb{Q}\),
\(\omega=\omega_{03}\,\Omega_0\).
If \(\omega_{03}=0\), then \(\omega=0\), so \(\det\omega=0\), contradicting nondegeneracy.
Thus \(\lambda=\omega_{03}\in\mathbb{Q}^{\times}\).
That is the only new step.
\end{proof}

\section{What is not proved}
The hypotheses are assumptions of the implication.
This note does not prove a functorial alternating form on a family, monodromy invariance as a geometric theorem, or nondegeneracy from a construction.
It does not prove holomorphic periods, existence of a compactification \(X\), \(\pi_1(X)=1\), or \(X\simeq S^6\).
The monodromy here is the rank-\(4\) group generated by \(T_1,T_2,T_0\), not sextic-fourfold monodromy.

\end{document}
